\section*{Capítulo \ref{ch:g5:decimals} --- Números decimais}

\begin{solution}{exo:g5:decimals:1}
$5.207$; \quad $0.85$; \quad $4.507$; \quad $12.009$.
\end{solution}

\begin{solution}{exo:g5:decimals:2}
O $4$ conta os décimos e o $8$, os milésimos:
\[
27.418 = 27 + \frac{4}{10} + \frac{1}{100} + \frac{8}{1000}.
\]
\end{solution}

\begin{solution}{exo:g5:decimals:3}
$3.60 = 3.6$: verdadeiro (um \omterm{def:g1:counting:zero}{zero} no fim não acrescenta nada).

$0.5 = 0.05$: falso --- $5$ décimos contra $5$ centésimos.

$2.070 = 2.07$: verdadeiro.

$1.3 = 1.03$: falso --- $3$ décimos contra $3$ centésimos.
\end{solution}

\begin{solution}{exo:g5:decimals:4}
$4.8 > 4.79$; \quad $0.302 < 0.32$ (décimos iguais, centésimos
$0 < 2$); \quad $6.5 = 6.500$; \quad $9.09 < 9.1$.
\end{solution}

\begin{solution}{exo:g5:decimals:5}
$3.041 < 3.104 < 3.14 < 3.4 < 3.41$.
\end{solution}

\begin{solution}{exo:g5:decimals:6}
$6.83$: entre $6$ e $7$; entre $6.8$ e $6.9$.

$0.492$: entre $0$ e $1$; entre $0.4$ e $0.5$.

$12.06$: entre $12$ e $13$; entre $12.0$ e $12.1$.
\end{solution}

\begin{solution}{exo:g5:decimals:7}
$8.276$: à unidade, $8$; ao décimo, $8.3$; ao centésimo, $8.28$. E
$3.95$ ao décimo: o algarismo dos centésimos é $5$, então sobe:
$4.0$.
\end{solution}

\begin{solution}{exo:g5:decimals:8}
Por exemplo $2.65$; \quad $0.415$; \quad $5.995$. (Há infinitas
respostas em cada caso.)
\end{solution}

\begin{solution}{exo:g5:decimals:9}
$0.85 < 0.9 < 1.02 < 1.2$ (em metros).
\end{solution}

\begin{solution}{exo:g5:decimals:10}
Compare primeiro os décimos: $7.12$ tem $1$ décimo e $7.9$ tem $9$
décimos, logo $7.12 < 7.9$. Lúcia comparou as partes decimais como
números inteiros ($12$ contra $9$), esquecendo que $12$ centésimos é
bem menos que $90$ centésimos.
\end{solution}

\begin{solution}{exo:g5:decimals:11}
Arredondando ao décimo, dá $6.4$ para os números de duas casas
decimais de $6.35$ (o meio arredonda para cima) até $6.44$: os
números $6.35$, $6.36$, \dots, $6.44$. Menor: $6.35$; maior: $6.44$.
\end{solution}
